\documentclass[11pt]{article}
\usepackage[paperwidth=210mm,paperheight=297mm,textwidth=182mm,top=20mm,bottom=22mm]{geometry}
\usepackage{amsmath,amssymb,newtxmath}
\usepackage[no-math]{fontspec}
\setmainfont{Times New Roman}
\usepackage{xeCJK}
\setCJKmainfont[AutoFakeBold=2]{SimSun}
\usepackage{graphicx,booktabs,array,hyperref,placeins}
\hypersetup{hidelinks,unicode=true}
\raggedbottom
\setlength{\emergencystretch}{2em}
\renewcommand{\arraystretch}{1.15}
\usepackage{seqsplit}
\begin{document}
\title{Litz 线分层建模与数值验证：完整研究报告 · 07 RT0 三维外场算例}
\author{缪雨轩\\淮南师范学院}
\date{}
\maketitle
\noindent 合并稿对应章节 / Chapter in merged manuscript: \href{report_zh-CN_single.pdf#chapter-07}{07}\par
\FloatBarrier
\section*{07 RT0 三维外场算例}
\hypertarget{chapter-07}{}
\pdfbookmark[0]{07 RT0 三维外场算例}{bookmark-chapter-07}
RT0 算例固定为 1×1×2 mm 均质铜长方体、200 kHz、Bz=1 mT RMS 外场和零净输运。三分量电流通过面通量自由度表示，连续性、公共端条件与外场功率分别检查。统一 q14 的 8192、16000、27648 单元损耗为 0.927463、0.933539、0.936927 mW；末两次变化为 0.655090\% 和 0.362916\%，均低于 1\%。末档 q18 对照的损耗相对变化为 1.54941×10⁻⁸，复面通量变化为 1.39291×10⁻⁸，低于 0.1\% 积分门槛。


场图由保存的面通量恢复。中心 J 采用相对两面的通量平均除以单元面积；单元平均损耗密度采用完整 RT0 局部重叠积分，不能用中心 |J|²/σ 替代。三维箭头仅显示固定相位零时刻的实部，按统一比例缩放；色标使用复 RMS 范数。切面的实际 z 位置、稀疏矢量抽样和整体积分分别标注，避免把抽样图看成完整场积分或全场最大值。


各档归档的 GMRES 迭代为 35 次，真残差满足 10⁻¹¹ 目标。原 q18 摘要中 null 指标及失败状态仍原样保留，新 CSV 后处理验收另列。旧致密算法的内存拒绝与新版矩阵自由成功求解是不同实现的证据，不能合并为一个状态。当前结论限于固定外场长方体，尚未覆盖曲线 Litz 掩膜、独立边界通量验证或实际输运端接。


\begin{figure}[!htbp]
\centering
\setcounter{figure}{12}
\refstepcounter{figure}\label{fig:rt0_sections}
\hypertarget{fig-rt0_sections}{}
\includegraphics[width=182mm]{../../../composites/zh-CN/rt0_sections.pdf}
\par\vspace{4pt}{\small\centering 图 13. RT0 外场电流与损耗切面。\par}
\vspace{3pt}{\small\raggedright 铜长方体为 1×1×2 mm，200 kHz、1 mT RMS 外场、零净输运。 (a) 单元平均损耗密度；(b) 复 RMS 电流范数。 电流由相对面通量恢复，单元平均损耗采用 RT0 重叠积分。切面位置和色标对应归档的重构定义。\par}
\end{figure}
\begin{figure}[!htbp]
\centering
\setcounter{figure}{13}
\refstepcounter{figure}\label{fig:rt0_vectors}
\hypertarget{fig-rt0_vectors}{}
\includegraphics[width=88.9mm]{../../../panels/zh-CN/rt0_vectors.pdf}
\par\vspace{4pt}{\small\centering 图 14. RT0 重构电流三维矢量。\par}
\vspace{3pt}{\small\raggedright 铜长方体为 1×1×2 mm，200 kHz、1 mT RMS 外场、零净输运。 箭头为固定相位的实部，色标为复 RMS 范数。 矢量采用固定比例与稀疏单元采样。单元与面通量身份在保存的重构数据中对应。\par}
\end{figure}
\begin{figure}[!htbp]
\centering
\setcounter{figure}{14}
\refstepcounter{figure}\label{fig:rt0_spatial}
\hypertarget{fig-rt0_spatial}{}
\includegraphics[width=88.9mm]{../../../panels/zh-CN/rt0_spatial.pdf}
\par\vspace{4pt}{\small\centering 图 15. RT0 空间与积分对照。\par}
\vspace{3pt}{\small\raggedright 铜长方体为 1×1×2 mm，200 kHz、1 mT RMS 外场、零净输运。 q14 空间序列与末档 q18 对照使用相同外场工况。 空间序列为 8192、16000、27648 单元。积分对照在最终成功网格上进行。\par}
\end{figure}
\end{document}
